\section{The Hebrew Verity: Making the Latin Bible}

\subsection{From the Seventy to the Hebrew}

The Bethlehem decades transformed Jerome's biblical program. In Rome he had
revised the Gospels against the Greek; in Palestine he first revised Old
Testament books against Origen's Hexaplaric Septuagint---the ``Gallican''
Psalter, made from that source, is the lasting monument of this stage---and
then, about 390, took the step that defined him: translating the Old
Testament directly from what he called the \textit{Hebraica veritas}, the
Hebrew verity, with the help of Jewish teachers. Between about 390 and 405/406
the whole Hebrew canon was done: Samuel and Kings first, then Psalms from the
Hebrew, Job, the Prophets, and onward to the Pentateuch and the remaining
books, each shipped with a preface answering critics in advance.

The preface to Samuel and Kings---the ``helmeted preface,'' as he armed
it---doubled as a manifesto on the canon. The Hebrews, he explains, count
twenty-two books, ``as, then, there are twenty-two elementary characters by
means of which we write in Hebrew all we say''; this preface ``may serve as a
`helmeted' introduction to all the books which we turn from Hebrew into
Latin,'' so that ``what is not found in our list must be placed amongst the
Apocryphal writings. Wisdom, therefore, which generally bears the name of
Solomon, and the book of Jesus, the Son of Sirach, and Judith, and Tobias,
and the Shepherd are not in the canon'' (Preface to Samuel and Kings). The
Church did not follow him here---the deuterocanonical books remained in her
Bible, and the Council of Trent would later define their status---but his
sharp distinction echoed through every later canon debate. His practice was
softer than his principle: pressed by bishops, he turned Tobit into Latin in
``the work of one day,'' dictating from an interpreter's Aramaic, and Judith
in a single sitting, noting that ``the Nicene Council is considered to have
counted this book among the number of sacred Scriptures'' (Prefaces to Tobit
and Judith). Wisdom, Sirach, Baruch, and Maccabees he declined to touch.

\subsection{What ``the Vulgate'' is and is not}

\witnesslimit{``Jerome translated the Vulgate'' needs three corrections, all
grounded in his own prefaces and the standard textual scholarship. First, in
the New Testament his prefaced, documented work is the revision of the four
Gospels; whether he ever revised the rest of the New Testament is disputed
(his own catalogue claims ``the New Testament from the Greek,'' but no
prefaces exist, and critical opinion has long been divided---the 1915
International Standard Bible Encyclopedia already weighed ``whether Jerome
revised the whole New Testament or only the Gospels''), and much modern
scholarship assigns the remaining books' revision to unknown hands. Second,
the Psalter that the Vulgate carried through history is not his Hebrew
version: ``His Psalter from the Hebrew never ousted the Gallican which still
holds its place in the Vulgate'' (ISBE). Third, the received Vulgate is a
composite: Jerome's Hebrew-based Old Testament and revised Gospels, his
hasty Tobit and Judith, and ``the unrevised Old Latin (Wisdom,
Ecclesiasticus, 1 and 2 Maccabees and Baruch)'' (ISBE). The name
``Vulgate'' itself, and the book's status as the Latin church's common
Bible, are later developments; Trent's decree and every subsequent official
edition concern that composite, not a single man's autograph.}

\subsection{The gourd at Oea: Augustine's challenge}

The Hebrew program had distinguished critics, and the greatest was Augustine
of Hippo. Their correspondence---begun awkwardly when early letters
miscarried and Jerome heard that a book against him was circulating---turned,
by 403, to the principle of the thing. Augustine ``would much rather'' Jerome
translated from the Septuagint, fearing that a Hebrew-based text would split
the Church's two lungs: ``it will be a grievous thing that, in the reading of
Scripture, differences must arise between the Latin Churches and the Greek
Churches.'' And he had a parable from life: at Oea in Tripolitania a bishop
had read Jerome's Jonah, and at the plant over Jonah's head---Jerome's
\textit{hedera}, ivy, where the old version had gourd---``arose such a tumult
in the congregation, especially among the Greeks, correcting what had been
read, and denouncing the translation as false,'' that the bishop consulted
the local Jews and restored the old reading rather than lose his people
(Augustine, \work{Letter} 71.4--5). Jerome's reply defended the philology
---the Hebrew word names a shrub with large leaves, ``ciceion,'' which he had
rendered by the nearest familiar plant---and returned the volley with aged
dignity: ``I beseech you to have some consideration for a soldier who is now
old and has long retired from active service, and not to force him to take
the field''; and, earlier in the exchange, ``do not, because you are young,
challenge a veteran in the field of Scripture\ldots\ The tired ox treads with
a firmer step'' (\work{Letters} 112.22, 102.2). Augustine, for his part, laid
down in this correspondence a rule the Church never forgot, that only
Scripture's own authors command unconditional assent: ``of these alone do I
most firmly believe that the authors were completely free from error''
(Augustine, \work{Letter} 82.3). The two doctors ended in peace and
collaboration against the Pelagians; the Church eventually kept both
verdicts, reading Jerome's Hebrew-based text while continuing to honor the
Septuagint the apostles had used.
